% ECONOMICS_PROBLEM: {"id": "C72-14", "title": "The Catch-Up optimal-draw conjecture for consecutive integers", "jel_code": "C72", "source_id": "OP-0589"}
% !TeX program = xelatex
\documentclass[11pt,a4paper]{article}
\usepackage[margin=23mm]{geometry}
\usepackage{fontspec}
\setmainfont{Latin Modern Roman}
\usepackage{amsmath,amssymb,mathtools}
\usepackage{xcolor,enumitem,booktabs,longtable,array,xurl,hyperref}
\definecolor{muted}{HTML}{53636C}
\hypersetup{colorlinks=true,urlcolor=blue,linkcolor=blue}
\setlength{\parindent}{0pt}
\setlength{\parskip}{5pt}
\newcommand{\R}{\mathbb R}
\newcommand{\E}{\mathbb E}
\newcommand{\N}{\mathbb N}
\newcommand{\ind}{\mathbf 1}
\newcommand{\Var}{\operatorname{Var}}
\newcommand{\Cov}{\operatorname{Cov}}
\newcommand{\supp}{\operatorname{supp}}
\DeclareMathOperator*{\argmin}{arg\,min}
\DeclareMathOperator*{\argmax}{arg\,max}
\newcommand{\entrymeta}[3]{{\sffamily\small\color{muted}\textbf{\texttt{#1}}\quad|\quad #2\par #3\par}}
\newcommand{\Problemref}[1]{\texttt{#1}}
\newcommand{\JELref}[2]{\texttt{#2} (JEL \texttt{#1})}
\begin{document}
\section*{The Catch-Up optimal-draw conjecture for consecutive integers}
\textbf{Problem ID:} \texttt{C72-14}\quad \textbf{JEL:} \texttt{C72}\par
% BEGIN ECONOMICS STATEMENT
\label{problem:OP-0589}
\label{supp:catchup}

\noindent\textbf{Original source:} \emph{Catch-Up on Consecutive Integers}, 8 October 2026.
\textbf{Formulation:} Unresolved conjecture in the supplied research note.

\subsection*{Definitions and assumptions}

Let $S$ be a finite set of distinct positive integers, initially with scores
$(s_1,s_2)=(0,0)$. In \emph{Catch-Up}, player 1 makes exactly one initial
selection. Thereafter the designated player removes and scores one remaining
integer at a time, continuing while strictly behind and passing the turn as
soon as their score equals or exceeds the opponent's score. The new player
must select at least one integer if any remain, even when the scores are tied.
If the active player cannot catch up, that player receives the remaining
integers and the game ends. At termination, compare the final scores. Treat
the final difference $s_1-s_2$ as the zero-sum payoff to player 1, so
\emph{optimal play} means minimax play for this exact scoring objective.

For a remaining set $R$ and deficit $d\ge0$ of the player next to select,
let $V(R,d)$ be that player's minimax final score advantage. The
one-selection game-value specification is
\[
V(\varnothing,d)=-d,\qquad
V(R,d)=\max_{x\in R}
\begin{cases}
	V(R\setminus\{x\},d-x),&x<d,\\[2pt]
	-V(R\setminus\{x\},x-d),&x\ge d,
\end{cases}\quad(R\ne\varnothing).
\]
Equality $x=d$ transfers the turn. For consecutive integers, write
$S_N=\{1,\ldots,N\}$ and $T_N=N(N+1)/2$.

\subsection*{Formal research question}

\textbf{Optimal-draw conjecture.} Determine whether
\[
\boxed{\quad \forall N\ge1,\quad
	N\equiv0\ \text{or}\ 3\pmod4
	\ \Longrightarrow\ V(S_N,0)=0.\quad}
\]
Equivalently, for every $N$ in these congruence classes, the minimax
outcome is an exact split of the available sum:
\[
s_1=s_2=\frac{T_N}{2}=\frac{N(N+1)}4.
\]
An equivalent opening-value formulation asks, for every such $N$, that
\[
V(S_N\setminus\{x\},x)\ge0\quad\text{for every }x\in S_N,
\qquad
\min_{x\in S_N}V(S_N\setminus\{x\},x)=0.
\]
Either prove this uniform assertion or exhibit a specified $N$ with
$N\equiv0,3\pmod4$ and $V(S_N,0)\ne0$.

\subsection*{Interpretation and scope}

Evenness of $T_N$ is a necessary arithmetic condition for a draw but is
not an equilibrium argument. Existence of a legal equal-sum allocation or
of some drawn sequence of moves differs from the minimax assertion.
The conjecture concerns all legal multiple-selection catch-up turns, not
only strategies that force opponents to stop after one selection. A
proposed induction or block strategy must account for the current score
deficit as well as the remaining set.

\subsection*{Source audit and status}

The supplied working note formulates the conjecture and reports partial
structural and computational evidence; it does not prove the assertion
for all $N$. The note attributes an opening symmetry to Bernhard von
Stengel and distinguishes finite computations from a general result.
The claimed computational checks in that note are not independently
verified for this catalogue. No proof, attempted invariant, or example
from the note is reproduced here.

\subsection*{References}

\begin{enumerate}
	\item[S1] A. Isaksen, M. Ismail, S. J. Brams, and A. Nealen (2015).
	\emph{Catch-Up: A Game in Which the Lead Alternates}.
	Game \& Puzzle Design 1(2), 38--49.
	\url{https://mpra.ub.uni-muenchen.de/108784/}.
	\item[S2] B. von Stengel (2019). \emph{Computational progress on the Catch-Up game}.
	CGTC III conference abstract. \url{https://ludicum.org/en/cgtc-iii/}.
	\item[S3] \emph{The Catch-Up Game Conjecture}, formal-conjectures issue 1324.
	\url{https://github.com/google-deepmind/formal-conjectures/issues/1324}.
	\item[S4] Supplied source: \emph{Catch-Up on Consecutive Integers: Precise
		formulation, rigorous partial results, proof attempts, and a counterexample
		to a proposed strategy}, 8 October 2026 (research working note).
\end{enumerate}

\subsection*{Common conventions: Source mathematical conventions}

\textbf{Finite games.} For a finite set $A$, $\Delta(A)$ is its probability
simplex. Players choose independent mixed strategies in Nash equilibrium;
correlation is allowed only when explicitly part of the solution concept.
In a game with payoff functions $u_i$ and strategy profile $x$, an additive
$\varepsilon$ Nash equilibrium satisfies
\[
 u_i(x)\geq u_i(x_i',x_{-i})-\varepsilon
 \quad\text{for every player $i$ and every admissible deviation $x_i'$}.
\]
For numerical approximation, payoffs are normalized as locally specified.
An $\varepsilon$ payoff residual is distinct from distance at most
$\varepsilon$ to the set of exact equilibria. Point convergence, convergence
of empirical play, and convergence in distance to an equilibrium set are
also distinct.

\textbf{Computational representations.} Unless stated otherwise, a complexity
question takes rational data with binary encoding; polynomial time is measured
in total bit length. A fixed-accuracy polynomial algorithm is not necessarily
polynomial in $1/\varepsilon$, and neither is necessarily polynomial in
$\log(1/\varepsilon)$. Succinct games and value, demand, or sampling oracles
must declare their interfaces. Exact real solutions require an explicit
representation; an unspecified list of arbitrary real numbers is not a
finite computational output. A claimed algorithmic barrier is meaningful
only for its specified model and reductions.

\textbf{Dynamic games.} Histories, information, observation, and allowable
randomization are part of the model. A stationary strategy depends only on
the current state; a positional strategy in a deterministic graph game
chooses one action at each controlled vertex. Nash equilibrium at an initial
state and equilibrium after every history are different requirements.
An expectation of a pathwise $\liminf$ payoff, a $\liminf$ of expected averages,
a limiting expected average, and a uniform finite-horizon equilibrium are
different criteria. None is silently substituted for another.

\textbf{Allocation and mechanisms.} A complete allocation assigns every item
exactly once unless otherwise stated. Zero-valued goods, negative values,
capacity constraints, feasibility, lotteries, and copies of goods affect
fairness definitions and are specified locally. Dominant-strategy incentive
compatibility, Bayesian incentive compatibility, universal truthfulness,
and truthfulness in expectation impose different inequalities. Whenever
an objective ratio has zero denominator, its convention is stated or those
instances are excluded.

\textbf{Infinite-dimensional models.} Expectations, integrals, stochastic
equations, and optimization objectives require their local regularity,
integrability, filtrations, and admissible domains. A backward--forward
mean-field system is a boundary-value problem; it is not an initial-value
semiflow without an additional construction. Long-time, large-population,
and vanishing-noise limits require an order of limits and a convergence
topology. If a minimum need not be attained, the objective is its infimum
or supremum and the approximation requirement is stated separately.
% END ECONOMICS STATEMENT
\end{document}
